\chapter{Numerical methods for the Vlasov--Maxwell--Proca system}\label{ch:numerics}

Chapter~\ref{ch:plasma} gave the analytic yardsticks; the numerical methods that must reproduce them come next. The paper compares four kinds of calculation: a particle-in-cell code, a semi-Lagrangian grid solver, a Hermite spectral solver in velocity, and the analytic theory of Chapters~\ref{ch:darkphoton} and~\ref{ch:plasma}. A reader who has used one family of methods tends to take for granted what the others do badly, and the conclusions of the paper depend on exactly those differences: a particle method has shot noise but no positivity problem, a spectral method has the reverse, and a grid method pays in memory for what it saves in noise. This chapter describes the methods in enough detail to make that statement precise, and then the architecture of Dark-SPECTRAX that implements the Hermite--Fourier one.

\section{Particle-in-cell}\label{sec:pic}

A particle-in-cell code represents the distribution function by a finite number of macroparticles, $f_s(x,v,t)\simeq\sum_pw_p\,\delta(x-x_p(t))\,\delta(v-v_p(t))$, each carrying a weight $w_p$ (number of physical particles) and moving along the characteristics of the Vlasov equation. The fields live on a grid. The method was developed in the 1960s and 1970s, and the standard references are the books by Birdsall and Langdon~\citep{BirdsallLangdon1985} and Hockney and Eastwood~\citep{HockneyEastwood1988} and the review by Dawson~\citep{Dawson1983}. A step has four parts: deposit charge and current from particles onto the grid, advance the fields, interpolate the fields back to the particles, and push the particles.

The grid enters through a shape function $S$. Each macroparticle is a finite-size cloud, and the charge density at node $x_j$ is $\rho_j=\sum_pq_pw_p\,S(x_j-x_p)/\Delta x$, with $\int S=1$. Using for $S$ the B-splines of order $m$ gives nearest-grid-point ($m=0$), cloud-in-cell ($m=1$), triangular-shaped cloud or quadratic ($m=2$) and the higher-order variants, up to the quintic splines ($m=5$) used by HHS. The field is interpolated to the particle with the same shape, $E(x_p)=\sum_jE_jS(x_j-x_p)$, because with identical deposition and interpolation the discrete force between two particles is antisymmetric and total momentum is conserved (the analogue of Newton's third law). Finite-size particles also change the dynamics: their Fourier transform $\hat S(k)$ multiplies the effective coupling, which reduces the strength of collisions at scales below the cloud size and damps the high-$k$ structure of the noise~\citep{LangdonBirdsall1970}. A higher-order shape is smoother and noisier at the grid scale only through the aliasing of $k$ components beyond the Nyquist wave number, so raising the order trades a larger stencil for a cleaner spectrum.

Maxwell's equations are advanced with a Yee-type staggered scheme, and the current must be deposited so that the discrete continuity equation holds, $\partial_t\rho_j+(J_{j+1/2}-J_{j-1/2})/\Delta x=0$, otherwise Gauss's law drifts. Villasenor and Buneman gave a charge-conserving deposition for first-order shapes~\citep{VillasenorBuneman1992}, and Esirkepov generalized it to arbitrary shape~\citep{Esirkepov2001}. In one dimension the construction is elementary: the current through a cell face is the change in the cumulative shape function that crosses the face during the step. A scheme that does not do this must correct the electric field with a Poisson solve.

The particle push for a relativistic particle uses $\bu=\gamma\bv$. The Boris algorithm~\citep{Boris1970} splits the Lorentz force $\dd\bu/\dd t=(q/m)(\bE+\bu\times\bB/\gamma)$ into two half-accelerations by the electric field and one rotation by the magnetic field:
\begin{align}
\bu^-&=\bu^{n}+\frac{q\Delta t}{2m}\bE ,\qquad \gamma^-=\sqrt{1+|\bu^-|^2/c^2},\\
\bu'&=\bu^-+\bu^-\times\mathbf t ,\qquad \bu^+=\bu^-+\bu'\times\frac{2\mathbf t}{1+t^2} ,\qquad\mathbf t=\frac{q\Delta t}{2m\gamma^-}\bB ,\\
\bu^{n+1}&=\bu^++\frac{q\Delta t}{2m}\bE .
\end{align}
The rotation preserves $|\bu|$ exactly, so the magnetic field does no work at any step size, and the map is volume preserving, which is why Boris is the default in codes that must run for very many steps. The rotation angle $\theta$ satisfies $\tan(\theta/2)=t$, a phase error of second order in $\Delta t$. In the relativistic limit Boris has a known flaw: in crossed $\bE$ and $\bB$ fields with $|\bE|\simeq c|\bB|$ the force on a particle that is moving at constant velocity does not cancel to the required accuracy, because the average of $\bv$ is taken at the half step with a $\gamma$ computed from the intermediate momentum. Vay proposed to evaluate the magnetic force with the average of the old and new velocities and solve the resulting implicit equation in closed form~\citep{Vay2008}. In his scheme one computes $\bu'=\bu^n+(q\Delta t/m)(\bE+\bv^n\times\bB/2)$ and then
\begin{equation}
\gamma_{n+1}^2=\tfrac12\Big(\sigma+\sqrt{\sigma^2+4(\tau^2+u_*^2)}\Big),\quad
\bu^{n+1}=\frac{\bu'+(\bu'\!\cdot\!\mathbf t)\mathbf t+\bu'\times\mathbf t}{1+t^2},
\end{equation}
with $\boldsymbol\tau=q\Delta t\bB/2m$, $\sigma=\gamma'^2-\tau^2$, $u_*=\bu'\cdot\boldsymbol\tau/c$ and $\mathbf t=\boldsymbol\tau/\gamma_{n+1}$. Both pushers reduce to the same thing in the electrostatic one-dimensional problem of this note, where $\bB=0$ and the push is the second-order leapfrog; the relevance of the choice to HHS is that they use the Vay algorithm, which we cannot vary in a code without a magnetic field.

The sampling error is the defining property of the method. If a cell of width $\Delta x$ holds $N_{\rm ppc}$ macroparticles, the density fluctuates by $\delta n/n\sim N_{\rm ppc}^{-1/2}$, so the PIC error converges as $N^{-1/2}$ unless the particles are loaded in a quiet start. The fluctuations drive electric fields: in electrostatics $E_k=\ii e\,\delta n_k/(\epsilon_0k)$, so the energy density of the noise in a cell is $\sim\tfrac12nT\,(\Delta x/\lDe)^2/N_{\rm ppc}$ (this estimate is ours, from $k\sim1/\Delta x$ and $\delta n^2/n^2=1/N_{\rm ppc}$). For the HHS set-up, $\Delta x=0.04\,c/\wpe$ and $\lDe=0.0316\,c/\wpe$, so the noise energy is about $1.6/N_{\rm ppc}$ times $\tfrac12nT=\Eth$; with the ladder of 103 to 1648 particles per cell of the paper this is $1.6\times10^{-2}$ to $10^{-3}$ of $\Eth$, small compared with the heating reported (tens of $\Eth$). The noise is nevertheless not innocuous. It seeds every unstable mode at a finite amplitude, which a deterministic code does not have unless it is given one, and the paper's central open question is whether this seed, and not a deterministic mechanism, produces the heating (Sec. ``Continuum methods against the PIC heating'' of the paper). The PIC has no positivity problem: the distribution is a sum of delta functions weighted positively.

Energy conservation in explicit PIC is only approximate. The scheme conserves momentum exactly, with matched shapes, but the energy changes by an amount that is second-order in $\Delta t$ and depends on the grid: when $\Delta x$ is larger than a few Debye lengths, aliased modes heat the plasma secularly, the finite-grid instability analyzed by Langdon and Birdsall~\citep{LangdonBirdsall1970}. Implicit schemes can conserve energy and charge simultaneously at the price of a nonlinear solve per step, as shown by Chen, Chac\'on and Barnes~\citep{ChenChacon2011} and by Markidis and Lapenta~\citep{MarkidisLapenta2011}. HHS use SHARP, a relativistic PIC code with a spatially higher-order discretization~\citep{Shalaby2017}, with fifth-order spline functions in deposition and interpolation, the Vay pusher, and between $2\times10^2$ and $2\times10^5$ particles per cell (Supplemental Material, Sec.~II of Ref.~\citep{HHS2026}). The independent relativistic PIC of the paper is built on Dark-JAX-in-Cell~\citep{DarkJIC}, with quadratic shapes by default and a quintic test. At $\vq=0.03$ the paper's electron heating at $\wpe t=1000$ converges to a plateau of about 25 initial thermal energies, with changes of 1.4\% and 0.1\% for each doubling of particles from 412 to 1648 per cell and 0.8\% when the cell is halved from $0.02$ to $0.01\,c/\wpe$ (paper, Table on PIC convergence). Dark-JAX-in-Cell was also used to check the Landau damping of the Hermite model: at 128 cells and 320\,000 particles the fitted frequency and damping agree within 0.1\% and 1\% (README, comparison with Dark-JAX-in-Cell), a measure of what the PIC needs to match a noise-free solver at $k\lDe=0.5$.

\section{Continuum methods: grids and discontinuous Galerkin}\label{sec:grid}

A continuum method represents $f$ on a grid in phase space, or by an expansion in basis functions, and has no sampling noise. Its costs are the curse of dimensionality (a six-dimensional grid for the full problem) and the need to keep $f$ positive and well resolved when it develops fine structure.

The semi-Lagrangian approach advances $f$ along characteristics over a time step and interpolates back to the grid. For the Vlasov--Poisson system the operator splitting of Cheng and Knorr~\citep{ChengKnorr1976} alternates a shift in $x$ by $v\Delta t$ and a shift in $v$ by $E\Delta t$; each is an exact translation, which makes the scheme second order, unconditionally stable in the translation steps (no CFL condition) and cheap, and with a spectral interpolation each shift is exact for band-limited data. The comparison of Eulerian solvers by Filbet and Sonnendr\"ucker~\citep{FilbetSonnendrucker2003} remains a good guide to the properties of this class, and Einkemmer and Ostermann gave the convergence analysis of Strang splitting for Vlasov-type equations~\citep{Einkemmer2014}. For the Vlasov--Maxwell system, Crouseilles, Einkemmer and Faou proposed a Hamiltonian splitting that handles the magnetic terms with a sequence of sub-steps that are each exactly solvable~\citep{Crouseilles2015}. The repository uses a NumPy semi-Lagrangian Vlasov--Poisson solver, extended to Vlasov--Amp\`ere--Proca for the dark runs, as an independent reference that shares no code with SPECTRAX (README, independent grid codes), and a relativistic momentum-space grid solver of the same family for the relativistic checks of the paper. In the relativistic grid, the momentum grid has 4096 nodes and $N_x=32$ Fourier modes (paper, Sec. on the relativistic extensions).

Discontinuous Galerkin methods approximate $f$ element by element with polynomials of fixed degree and couple elements by numerical fluxes. They are high order, local and conservative, and have been developed for kinetic equations in a series of works by Juno, Hakim and collaborators~\citep{Juno2018}; for the relativistic Vlasov--Maxwell system, Juno et al.\ recently built a conservative DG scheme with a velocity-space mapping that deals with the wide range of energy scales in relativistic plasmas, and describe its noise-free character as the point of the method~\citep{Juno2026}. This is the nearest neighbour to our relativistic grid and to the relativistic Hermite option in the sense that it addresses the same physics (relativistic kinetics without Poisson noise), but over a different basis.

Table~\ref{tab:methods} summarizes the properties that decide the choice for the problem at hand. The entries are qualitative and refer to the schemes described above, in their standard form; each can be modified, and the references show how.

\begin{table}[t]
\centering
\caption{Qualitative comparison of the methods in the standard forms described in this chapter. ``Noise'' is sampling noise in the distribution function. ``Positivity'' is whether $f\ge0$ is guaranteed by the discretization.}\label{tab:methods}
\small
\begin{tabular}{@{}p{3.2cm}p{1.6cm}p{2.3cm}p{3.5cm}p{3.4cm}@{}}\toprule
Method & Noise & Positivity & Invariants (semi-discrete) & Main cost driver\\\midrule
PIC, explicit & yes, $N_{\rm ppc}^{-1/2}$ & yes (weights $>0$) & momentum; energy only to $O(\Delta t^2)$, grid heating & particles per cell, steps\\
PIC, implicit energy-conserving & yes & yes & charge and energy & nonlinear solve per step\\
Semi-Lagrangian grid & no & not guaranteed (interpolation) & mass; energy approximately & phase-space grid points\\
DG in phase space & no & not guaranteed (limiters help) & mass, energy with central fluxes & elements $\times$ degree$^{d}$\\
Hermite, AW & no & no & mass, momentum, and energy up to closure; AW Galerkin instability & Hermite order, step $\propto1/\sqrt{N_n}$\\
Hermite, SW & no & only in the square-root form & $L^2$ norm; no closure conserves all three moments & Hermite order\\\bottomrule
\end{tabular}
\end{table}

\section{Spectral discretization in velocity: the Hermite expansion}\label{sec:hermite}

The velocity space of a plasma close to a Maxwellian is well represented by a short expansion in Hermite functions, and a spectral method turns the Vlasov equation into a hierarchy of equations for the expansion coefficients, which are moments of the distribution. Grad's moment method of 1949 is the antecedent~\citep{Grad1949}, Fourier--Hermite expansions were applied to the Vlasov equation from the 1960s, and Holloway's analysis of spectral velocity discretizations~\citep{Holloway1996} and the velocity-scaled Hermite simulations of Schumer and Holloway~\citep{SchumerHolloway1998} gave the theory and practice that most of the modern work follows. The modern revival, driven by exascale hardware and by the demand for noise-free kinetic simulation, includes the fully implicit scheme of Delzanno with exact discrete conservation~\citep{Delzanno2015}, the comparison of Camporeale et al.\ with PIC~\citep{Camporeale2016}, the three-dimensional Hermite--DG code of Koshkarov et al.~\citep{Koshkarov2021} and the analysis of its conservation properties and time integrators by Pagliantini et al.~\citep{Pagliantini2023a}.

We write the expansion in one velocity dimension; the three-dimensional case is a tensor product. Let $\xi=(v-u)/a$, with width $a$ and shift $u$ to be chosen. The \emph{asymmetrically weighted} (AW) basis, used by the parent solver SPECTRAX, is
\begin{equation}
f(x,v,t)=\sum_{n=0}^{N-1}C_n(x,t)\,\Psi_n(\xi),\qquad \Psi_n(\xi)=\frac{H_n(\xi)\,\ee^{-\xi^2}}{\sqrt{\pi\,2^n\,n!}} ,\label{eq:AW}
\end{equation}
where $H_n$ are the physicists' Hermite polynomials. The functions are orthogonal with respect to the weight $\ee^{+\xi^2}$ and not $\ee^{-\xi^2}$, hence the name: the Gaussian factor appears once in the basis function and the weight is the inverse Gaussian. The coefficients are moments: since $\int\Psi_nH_m\,\dd\xi=\sqrt{2^nn!}\,\delta_{nm}$,
\begin{equation}
C_n=\frac{1}{a\sqrt{2^nn!}}\int f\,H_n\!\Big(\frac{v-u}{a}\Big)\dd v .
\end{equation}
The density is $\int f\,\dd v=aC_0$, the particle flux is $\int vf\,\dd v=a\,(uC_0+aC_1/\sqrt2)$, and the second moment involves $C_2$. Number, momentum and energy are therefore the first three coefficients, and the current that sources Amp\`ere's law is a linear function of $C_1$ (and $C_0$ in a shifted basis), which is why the field solve costs nothing.

The Hermite polynomials satisfy $\xi H_n=\tfrac12H_{n+1}+nH_{n-1}$ and $H_n'=2nH_{n-1}$; from them, the ladder relations of the AW basis are
\begin{equation}
\xi\Psi_n=\sqrt{\tfrac{n+1}{2}}\,\Psi_{n+1}+\sqrt{\tfrac n2}\,\Psi_{n-1},\qquad \partial_\xi\Psi_n=-\sqrt{2(n+1)}\,\Psi_{n+1} .\label{eq:ladder}
\end{equation}
The first holds because $\xi$ multiplies $H_n$ and the normalization $1/\sqrt{\pi2^nn!}$ turns $\tfrac12H_{n+1}$ into $\sqrt{(n+1)/2}$ times the normalized $\Psi_{n+1}$; the second follows from $\partial_\xi(H_n\ee^{-\xi^2})=(2nH_{n-1}-2\xi H_n)\ee^{-\xi^2}=-H_{n+1}\ee^{-\xi^2}$. Insert $v=u+a\xi$ in the one-dimensional Vlasov equation, project on the basis and use Eq.~\eqref{eq:ladder}. In the Fourier representation in $x$ the result is
\begin{equation}
\partial_tC_n+\ii k\Big[u\,C_n+a\sqrt{\tfrac{n+1}{2}}\,C_{n+1}+a\sqrt{\tfrac n2}\,C_{n-1}\Big]=\frac{q}{ma}\sqrt{2n}\;\big(E\ast C_{n-1}\big) ,\label{eq:hierarchy}
\end{equation}
with $\ast$ the convolution in $k$ that represents the product $E(x)C_{n-1}(x)$. The streaming term couples neighbouring orders through the symmetric tridiagonal matrix $T$ with entries $T_{n,n\pm1}=\sqrt{(n+1)/2}$ or $\sqrt{n/2}$; its eigenvalues are the nodes of the Gauss--Hermite quadrature, the largest being about $\sqrt{2N}$. The force term only moves information from lower to higher $n$: it is strictly lower triangular in the Hermite index. This structure explains most of the numerical behaviour of the method.

Consider free streaming from a basis-matched Maxwellian, $f_1=\ee^{-\ii kvt}f_M(v)$ with $a=\sqrt2\vte$. Then $C_n=\int\ee^{-\ii s\xi}H_n(\xi)\ee^{-\xi^2}\dd\xi/(\sqrt\pi\sqrt{2^nn!})$ with $s=kat$, and the identity $\int H_n\ee^{-\xi^2-\ii s\xi}\dd\xi=\sqrt\pi(-\ii s)^n\ee^{-s^2/4}$ gives $|C_n|^2=(k\vte t)^{2n}\ee^{-(k\vte t)^2}/n!$. The spectrum in $n$ is a Poisson distribution with mean $(k\vte t)^2$: the information flows up the ladder at the rate with which the velocity structure $\ee^{-\ii kvt}$ becomes finer, and a truncation at $N$ is accurate until $(k\vte t)^2\sim N$, that is, $t\sim\sqrt N/(k\vte)$. After this time the truncated system reflects the wave at the end of the ladder and the solution recurs. This is the Hermite counterpart of the recurrence in a velocity grid with spacing $\Delta v$, which occurs at $t=2\pi/(k\Delta v)$. The fix in both cases is dissipation at the small scales, which in the Hermite case is a damping that acts on high $n$.

The closure problem is the question of what to set for $C_N$ in the equation for $C_{N-1}$. Truncation, $C_N=0$, reflects; the choices in use are filters, hypercollisions and Landau-fluid closures. Hammett and Perkins constructed a closure from the linear response that reproduces Landau damping with a few fluid moments~\citep{HammettPerkins1990}. Hypercollisions add $-\nu\,s(n)C_n$ with $s(n)\propto n^p$, sometimes with a floor so that the first conservation orders are not touched; the parent solver uses $s(n)\propto n(n-1)(n-2)$, which leaves $n=0,1,2$ (density, momentum, energy) unchanged. Parker and Dellar showed how the free energy flows up the ladder in linear Landau damping and back down in the nonlinear case, which is the physical picture behind a closure that must not reflect~\citep{ParkerDellar2015}. Camporeale et al.\ added an artificial collisional operator for the filamentation problem~\citep{Camporeale2016}, and Issan et al.\ compared artificial collisions, exponential filters and nonlocal closures in a systematic study of the effect on Landau damping, nonlinear Landau damping and the two-stream instability~\citep{Issan2025}. For the paper we use a hypercollision whose strength is scaled by the maximum electric field,
\begin{equation}
\nu_s=c\Big|\frac{q}{m}\Big|_s\frac{\sqrt{2N_n}\,\max|E-\langle E\rangle|}{a_s}\;\frac{n(n-1)(n-2)}{(N_n-1)(N_n-2)(N_n-3)},
\end{equation}
so that the rate at the top of the ladder equals $c$ times the Hermite-space velocity (in $n$ per unit time) at which the force carries information up the ladder (paper, Sec. on closures).

The alternative to the AW basis is the \emph{symmetrically weighted} (SW) basis $\varphi_n=H_n\ee^{-\xi^2/2}/\sqrt{\sqrt\pi\,2^nn!}$, orthonormal in $L^2$. Its ladder relations are $\xi\varphi_n=\sqrt{(n+1)/2}\,\varphi_{n+1}+\sqrt{n/2}\,\varphi_{n-1}$ and $\partial_\xi\varphi_n=\sqrt{n/2}\,\varphi_{n-1}-\sqrt{(n+1)/2}\,\varphi_{n+1}$, so that the discretized force term is an antisymmetric tridiagonal matrix times $E(x)$. A skew-symmetric Galerkin operator conserves the discrete $L^2$ norm of $f$ and cannot have exponentially growing modes. The price is that the SW expansion does not preserve the simple moment structure of AW, and its conservation of momentum and energy requires closures that Issan et al.\ show cannot simultaneously conserve the three invariants~\citep{Issan2024b}. The antisymmetric formulation, with the square-root transformation that also preserves positivity, is the subject of Ref.~\citep{Issan2024a}. Bessemoulin-Chatard and Filbet proved weighted-$L^2$ stability, and conservation of mass, momentum and energy, for the Hermite spectral form with a time-dependent weight~\citep{BessemoulinChatardFilbet2021}. A mixed Hermite--Legendre method that uses Legendre polynomials for localized features such as beams and plateaus is proposed in Ref.~\citep{Issan2026}, and a recent Galerkin method with AW functions that is designed to be stable is that of Dai~\citep{Dai2026}.

The distinction is not academic here, because the AW basis, which the parent solver uses and which makes the pump-frame remap exact, has a numerical instability. In a spatially varying field the frozen-field spectrum of the AW Galerkin operator $-\ii kaT+(q/ma)\sqrt{2n}\,E$ has modes that grow, concentrated in the top third of the ladder: the streaming matrix is real symmetric and the force matrix is strictly lower triangular and nilpotent, so neither has an eigenvalue with positive real part, but their sum is a non-normal matrix, and the product of upward and downward couplings can produce growth. At $K=5$ Fourier modes, $N_n=128$ and field amplitude $E_1=10^{-2}$ the largest growth rate is $1.6$ for the AW basis against $8\times10^{-15}$ (round-off) for the SW basis (\texttt{studies/highk/eig.json}, paper Fig.~2). The growth rate rises with $N_n$, which means that refining the velocity basis moves the end of the run earlier, and this defines the certified windows of the nonrelativistic paper: $1000$, $892$ and $791.5$ at $\vq=0.03$, $0.05$, $0.07$. The paper and the design note record the routes tried to remove the instability (stronger closure, GPU resolution, the SW basis and a square-root form) and why they failed under the declared gates; for the purpose of this chapter the point is that a Hermite basis can be energetically fine in a continuous sense and unstable in a discrete one, and the choice among bases is a choice of which defect to live with.

\section{Scaling, shifting and adaptivity}\label{sec:scaling}

The width $a$ and shift $u$ in Eq.~\eqref{eq:AW} are free parameters of the discretization. A basis that does not match the distribution needs many more terms, because the Hermite functions are orthogonal for one Gaussian only. Tang observed that the standard normalized Hermite functions require too many terms for Gaussian-type functions unless a scaling factor is chosen, and gave the factor in terms of the largest root of $H_{N+1}$~\citep{Tang1993}. Ma, Sun and Tang made the scaling depend on time to follow the spreading of the solution of a parabolic equation in an unbounded domain, and proved convergence~\citep{MaSunTang2005}. Schumer and Holloway introduced a variable velocity scale into the Hermite basis for the one-dimensional Vlasov--Poisson system and found errors reduced by orders of magnitude at no increase in workload~\citep{SchumerHolloway1998}; they also compared the two weightings and noted that the AW basis conserves particles and momentum exactly but not the square integral of $f$, and is numerically unstable, while the SW basis conserves the square integral (this is the instability that reappears in the next section and in the paper)~\citep{SchumerHolloway1998}. Fatone et al.\ studied the dependence on the scaling parameter in a semi-Lagrangian spectral Vlasov--Poisson scheme and concluded that the parameter should be updated during the computation~\citep{Fatone2018}. Pagliantini, Delzanno and Markidis designed the adaptivity: after each time step a criterion sets $a$ and $u$ from the bulk velocity and temperature of each species, and a map between the approximation spaces of different parameters preserves mass, momentum and energy; the resulting method has fluid--kinetic coupling, since the lowest coefficients become the fluid moments~\citep{Pagliantini2023b}.

The AW basis makes the remap simple. The coefficients with respect to a new pair $(u',a')$ are $C_m'=(a'\sqrt{2^mm!})^{-1}\int f\,H_m((v-u')/a')\,\dd v$, and $H_m((v-u')/a')$ is a polynomial of degree $m$ in $\xi=(v-u)/a$. Hence $C_m'$ is a linear combination of the moments $\int f\,\xi^j\,\dd v$ with $j\le m$, which are combinations of $C_0,\dots,C_j$: \emph{the remap is lower triangular}, $C'_m=\sum_{n\le m}R_{mn}C_n$, with no error for any $m<N$ when the lower coefficients are exact. In particular every moment below the truncation order is preserved exactly. The SW basis has no such property: the change of the Gaussian weight is not a polynomial map and mixes all orders. This is the reason why the model's ``pump frame'' uses AW functions: in the resonant calculation, each species oscillates with a coherent velocity $u_s(t)$ that grows to $0.5c$, and the centre follows the total uniform force, $\dot u_s=(q/m)_s(\bar{\bf E}+u_s\times\bar{\bf B})$ with the box averages including the dark field and the drive. This is an exact change of variables (the kernel receives $E-\bar E$), so the coherent $k=0$ quiver is carried by $u_s(t)$ and not by the high orders of the Hermite ladder, which would have to represent a displaced Gaussian by an enormous number of terms. Widths and residual drifts change only between integration segments, through the triangular remap, with the triggers and caps recorded in the repository ($|U-u|/a>0.3$ or $|a'/a-1|>0.05$; $|\Delta u|\le a'$; $a/a'\le1.1$).

\section{Fourier pseudo-spectral discretization in space}\label{sec:fourier}

In space the periodic box is discretized by a Fourier series, and products such as $E(x)C_{n-1}(x)$ in Eq.~\eqref{eq:hierarchy} are computed pseudo-spectrally: transform to the grid, multiply pointwise and transform back. The product of two fields with wave numbers $|m_1|,|m_2|\le M$ has wave numbers up to $2M$, and on a grid of $N_x$ points a wave number $p$ is indistinguishable from $p-N_x$. If all modes with $|m|\le M$ are retained, a product term at $p>N_x-M$ is aliased into a retained mode. Requiring that no product mode of order $2M$ alias into a retained mode, $2M<N_x-M$, gives $M<N_x/3$: the $2/3$ rule of Orszag~\citep{Orszag1971}, which discards the upper third of the spectrum and makes quadratic products exact. The parent applies a mask retaining modes with $3|m|<N_x$ (a ``strict'' $2/3$ rule, after SPECTRAX pull request 55), so that every grid size, including multiples of 3, is free of aliasing for quadratic products. The dealiased method is spectrally accurate for smooth fields; for the Vlasov equation the discretization in $x$ is not the limiting factor in the cases of interest, since the high-$k$ content of the fields is small or physically relevant (the OTSI cascade reaches large $k$ and drives the choice of $N_x$ in the paper).

A real field requires only $k\ge0$ (the real FFT), and the box average is a weighted sum over this half spectrum with weight 1 for $k=0$ (and for the Nyquist mode when $N_x$ is even) and 2 otherwise. This is how the energy ledger is evaluated from the stored coefficients without returning to the grid.

The two discretizations together set the cost of a run, and the numbers can be estimated before running. In the HHS box, $L=40\,c/\wpe$ and $\lDe=\vte/\wpe=0.0316\,c/\wpe$, so a wave number $k\lDe=1$ is $k=31.6\,\wpe/c$, which is the Fourier mode $m=kL/2\pi\simeq201$. With the $2/3$ rule the grid must have at least $3m\simeq604$ points to resolve that mode without aliasing, so $N_x=32$, $64$ and $128$ retain modes up to $m=10$, $21$ and $42$ and cover $k\lDe$ up to about $0.05$, $0.10$ and $0.21$, while $N_x=1024$ reaches $m=341$ and $k\lDe\simeq1.7$. The Hermite order needed for a collisionless linear phase-mixing problem up to the time $T$ is $N_n\gtrsim(k\vte T)^2$ from the Poisson spectrum of Eq.~\eqref{eq:hierarchy} (for $k\lDe=0.3$ and $\wpe T=100$ this is $\simeq900$, the order of the $N_n=512$ to $1024$ used for nonlinear Landau damping in the repository, and for $k\lDe=1$ and $\wpe T=1000$ it is $10^6$, which is why closures and collisions are needed and why the solution is only trusted until the first disagreement between orders). The cost of a PIC run scales as $N_xN_{\rm ppc}N_t$ with a noise error $\propto N_{\rm ppc}^{-1/2}$, while the Hermite run scales as $N_xN_nN_t$ with an error that decays exponentially in $N_n$ for smooth data until the filamentation reaches the end of the ladder. The comparison of Camporeale et al.\ shows the consequence: for the smooth problems of their study the Hermite method reached solutions several orders of magnitude more accurate at a fraction of the cost of PIC~\citep{Camporeale2016}. Our problem differs by being strongly nonlinear with an instability cascade to small scales, so that the Hermite advantage is not a given; the certified windows of the paper are the measure of how far it holds.

\section{Time integration}\label{sec:time}

The semi-discrete system is a large system of ordinary differential equations $\dot y=f(t,y)$ with a stiff and a nonstiff part. Streaming has imaginary eigenvalues $\pm\ii kax_j$, with $x_j$ the Gauss--Hermite nodes up to $\sqrt{2N_n}$, so the stable step of an explicit method decreases as $1/(K\sqrt{N_n})$. This is the CFL condition of the spectral method, and it is the cost of high resolution in $n$. The closure damping is real and negative with magnitude up to its strength, and is stiff only for strong closures.

The default integrator of Dark-SPECTRAX is the eighth-order explicit Runge--Kutta method Dopri8 of the Diffrax library~\citep{Kidger2021}, with embedded error estimation and a proportional-integral-derivative step-size controller that adjusts the step to a relative tolerance $10^{-10}$ and an absolute tolerance $10^{-12}$. A high order is appropriate because the problem is smooth, the cost per step is eight right-hand-side evaluations with a single FFT of the fields, and it is the integrator error and not the spatial or velocity truncation that we want to be negligible. The work ledger is integrated as additional variables with the same stages, so that $K-K_0=W_\gamma+W_D+\Wext$ is a statement about the integrator and closes to the solver tolerance.

The implicit midpoint rule is a different tool for a different purpose. For $\dot y=f(y)$ it reads $y_{n+1}=y_n+\Delta t\,f\big(\tfrac12(y_n+y_{n+1})\big)$. It is second-order, A-stable and symmetric, and it conserves every quadratic invariant of the continuous system. To see this, suppose $y^{\rm T}Qy$ is invariant with $Q$ symmetric, so that $y^{\rm T}Qf(y)=0$ for all $y$. Then $y_{n+1}^{\rm T}Qy_{n+1}-y_n^{\rm T}Qy_n=(y_{n+1}-y_n)^{\rm T}Q(y_{n+1}+y_n)=2\Delta t\,f(\bar y)^{\rm T}Q\bar y=0$ with $\bar y=\tfrac12(y_n+y_{n+1})$. For the Vlasov--Poisson and Vlasov--Maxwell systems in the Hermite representation the discrete energy and, in the SW basis, the $L^2$ norm are quadratic, so the midpoint rule conserves them exactly, which is the reason why implicit second-order schemes of this type (Crank--Nicolson, which is the midpoint rule for these equations) are used by Delzanno~\citep{Delzanno2015}, Issan et al.~\citep{Issan2024a} and Pagliantini et al.~\citep{Pagliantini2023b}. The parent solver provides it as a Diffrax solver with a Newton--GMRES nonlinear solve, and the cost is that of a nonlinear solve per step, not eight right-hand sides. Pagliantini et al.\ analyze when explicit methods, with relaxation, can preserve energy instead~\citep{Pagliantini2023a}.

When the stiffness comes from a known linear part the exponential integrators remove it from the stability restriction. Write $\dot y=Ay+g(t,y)$ with $A$ the stiff operator. Lawson's integrating-factor method sets $z=\ee^{-tA}y$, so that $\dot z=\ee^{-tA}g(t,\ee^{tA}z)$, and advances $z$ with an explicit Runge--Kutta method; for the linear test problem the stiffness is removed exactly and the explicit method sees only $g$~\citep{Lawson1967}. Cox and Matthews gave exponential time differencing as the other principal family~\citep{CoxMatthews2002}, and Hochbruck and Ostermann reviewed both, including the rate of convergence for semilinear problems~\citep{Hochbruck2010}. The review makes a point that matters here: for stiff problems the classical order is not always attained, and the Lawson--Runge--Kutta schemes are of lower stiff order than their nonstiff order. The effective order depends on the ratio of the step to the stiff time scale, and one should not claim the full order of Dopri8 when the closure rate times the step is large. For kinetic equations, Crouseilles, Einkemmer and Massot studied exponential integrators for collisionless kinetic equations, whose linear part (free streaming) is exactly solvable in Fourier space and removes the CFL condition. They showed that classic exponential integrators behave erratically in the step size they allow, and that Lawson methods do not suffer from this, which they demonstrated on the Vlasov--Poisson system and a drift-kinetic model~\citep{Crouseilles2020}.

In Dark-SPECTRAX the stiff part is the diagonal closure, $A=-\mathrm{diag}(r)$, which is exactly exponentiable. For $c=64$ the Jacobian's spectral radius equals the closure rate to the printed precision, and the Lawson wrapper around Dopri8 reduces the number of steps by a factor of $5.0$ to $\wpe t=500$ with $\Delta K_e$ and $\Wext$ agreeing with the plain integration to $5\times10^{-10}$; for $c=4$ the closure is not stiff and the step ratio is $0.96$ (\texttt{studies/expint/summary.json}). An exact streaming factor in the same framework gave no step reduction and cost 2.7 times more per step. The gain is real but it is not the limit of the method: the largest real part of the eigenvalues of the instability is 2--3, smaller than the closure rate, so the time step is not controlled by the instability (paper, failed-checks table).

\section{Architecture of Dark-SPECTRAX}\label{sec:architecture}

Dark-SPECTRAX is a Python package of about a thousand lines that adds the Proca sector, the prescribed drive, the work ledger, the pump-frame logic, a relativistic option and the study tooling to the Hermite--Fourier kernel of the parent solver SPECTRAX~\citep{SPECTRAX}. The parent is a JAX code, so that the right-hand side is a pure function of a state pytree, can be compiled with XLA, run on CPU or GPU, and differentiated. Dark-SPECTRAX keeps the parent as a dependency, pinned to a commit on an integration branch (the commit is recorded in each study record), and calls it through a small adapter section of one module; its own files are the model (configuration and right-hand side), the simulation drivers, the closure and the streaming factor, and the diagnostics.

The state is a dictionary of complex arrays. The Hermite coefficients are stored as \texttt{Ck} with shape $(N_sN_nN_mN_p,\;N_y,\;N_x/2+1,\;N_z)$: the first axis is the flattened species and Hermite index, $N_n,N_m,N_p$ being the numbers of Hermite orders in the three velocity directions, with the parent storage order $(p,m,n)$ in the Hermite index; the spatial axes are the real-FFT layout, with the half spectrum along $x$. The ordinary fields are \texttt{Fk} of shape $(6,\dots)$ for $\bE$ and $\bB$ in the same Fourier layout; the dark fields are \texttt{Dk} of shape $(10,\dots)$ for $\ED$, $\BD$, $\AD$ and $\phiD$; \texttt{W} holds the three works; and in the pump frame \texttt{B} holds the centres and widths $(u_s,a_s)$ as variables that evolve. All state is complex because the parent keeps its Fourier coefficients complex; the ledger variables are stored as complex with zero imaginary part because Diffrax requires one dtype across the Runge--Kutta stage buffers.

The right-hand side \texttt{rhs(t, y, model)} has one structure for all modes. The physical force fields are synthesized from \texttt{Fk} (plus $\eta\,$\texttt{Dk} in the self-consistent mode, plus the drive in the prescribed mode), multiplied by the dealiasing mask and transformed to the grid; one call to the parent's \texttt{Hermite\_Fourier\_system} applies streaming, the Lorentz force and the numerical hypercollision to \texttt{Ck}; the optional Dougherty collision operator is added from the parent (the multispecies forms of Sec.~\ref{sec:collisions}); the plasma current $\bJ$ is computed once and feeds both Amp\`ere laws, the ordinary one with $-\bJ/\Omega_{c0}$ and the dark one with $-\eta\bJ/\Omega_{c0}$, in the parent normalization; the field updates are the curl terms of Sec.~\ref{sec:coupling} in Fourier space. The powers $\Omega_0\avg{\bJ\cdot\bE}$, $\Omega_0\eta\avg{\bJ\cdot\ED}$ and $\Omega_0\avg{\bar{\bJ}\cdot\Edrive}$ are evaluated from the half-spectrum coefficients with the weights of Sec.~\ref{sec:fourier} and form \texttt{dW}. The parent normalization has velocities in units of $c$ and requires $m_s\Omega_{cs}=\Omega_0$ for the ledger to close; the model defines the species masses as $m_s=\Omega_0/\Omega_{cs}$ for this reason. The \texttt{inv\_c2} parameter (zero for the nonrelativistic path, which is then bitwise the old code) turns on the relativistic option, where the Hermite functions are in momentum.

The drivers sit on top of Diffrax. \texttt{run} compiles an ahead-of-time solve once and calls it, with Dopri8, a PID controller and an optional floor on the error norm of the high Hermite rows (a ``noise floor'' that stops the controller from chasing round-off in unseeded rows); the optional \texttt{integrating\_factor} wraps the solver in the Lawson solver, using the diagonal closure rate from the closure module or the exact streaming factor of the streaming module. \texttt{run\_adaptive} advances in segments between which \texttt{adapt\_basis} applies the triangular remap, either live from the trigger or replaying a recorded schedule, so that every event is recorded and reproducible. The state is saved with its Git commit and a run record, which is how the studies of the repository attach each number of the paper to a file. The solver interface is deliberately small: a model dataclass, an initial state built by \texttt{maxwellian} (basis-matched drifting Maxwellians with density perturbations), \texttt{consistent\_fields} (fields in Gauss equilibrium with the charge density, including the Yukawa near field of the dark photon) or \texttt{proca\_mode} (a free vacuum Proca mode), and the drivers.

The code is differentiable. \texttt{solve} returns only JAX arrays, so it composes with \texttt{jit}, \texttt{grad} and \texttt{jvp}; \texttt{with\_params} makes $\eta$, $\OmD$, the widths and the centres traced values. The adjoint is Diffrax's \texttt{RecursiveCheckpointAdjoint}, which stores a logarithmic number of checkpoints and recomputes in the backward pass, or forward mode for \texttt{jacfwd}. The test suite compares reverse-mode, forward-mode and adaptive-step gradients with centred finite differences (agreement of $10^{-10}$ between the two AD modes and $10^{-6}$ with finite differences in the test of the dark work with respect to $\eta$ and $\OmD$). A limitation is recorded in the docstring: with a fixed step and many saves on step boundaries, the Diffrax checkpointing was measured wrong by a few per cent in the version used, and the repository recommends the direct or forward adjoint there. Gradients through the pump-frame remap and through nonlinear-regime objectives are not validated. Differentiability matters for this problem in two ways. It allows inference of $(\eta,\OmD)$ from a measured energy deposition, which is the sort of use a constraint calculation would make of a noise-free solver, and it gives the sensitivity of the transfer to parameters without the sampling noise that would swamp a finite difference in a PIC code.

What this architecture does not do is also part of the description. The velocity space of the production runs is one- or two-dimensional in the Hermite index ($N_m,N_p$ small), the space is one-dimensional, the pusher is not a particle pusher, and the Hermite option for relativistic dynamics has no collisions. These are the limits that the paper lists as not yet established, and they bound what the continuum method can say about the early-Universe plasma.
